% F21：本书原创矢量示意；依据所在节的定义、证明或明确模型。
\begin{figure}[htbp]
\centering
\begin{tikzpicture}[x=1cm,y=1cm,>=stealth,every node/.style={font=\small},line width=.65pt]

\draw[->,AnalysisBlue,thick](1,.5)--(3,.5)node[right]{\(e_1\)};
\draw[->,AnalysisBlue,thick](1,.5)--(1,2.5)node[above]{\(e_2\)};
\draw[->,orange!80!black](1.7,.5)arc(0:90:.7);
\node at(2,-.3){\(e_1\wedge e_2\)};
\draw[->,AnalysisBlue,thick](5,.5)--(5,2.5)node[above]{\(e_2\)};
\draw[->,AnalysisBlue,thick](5,.5)--(7,.5)node[right]{\(e_1\)};
\draw[->,orange!80!black](5,1.2)arc(90:0:.7);
\node at(6,-.3){\(e_2\wedge e_1=-e_1\wedge e_2\)};
\draw[->,thick](9,0)--(12,0);\draw[->,thick](12,0)--(10.5,2.6);\draw[->,thick](10.5,2.6)--(9,0);
\draw[->,orange!80!black](10.5,0)--(10.5,-.7)node[below]{外法向};
\node at(10.5,3.3){外法向放首位};

\end{tikzpicture}
\caption[定向、换基与边界方向]{定向、换基与边界方向。交换有向基的次序改变外积符号；右图采用平面的标准定向，边界逆时针，外法向在有序基中放首位。}
\label{b1:fig:visual-f21}
\end{figure}
